\section{Thiết lập thực nghiệm}
Để xác định hiệu quả của Kiến trúc Ensemble lai đề xuất kết hợp với chiến lược Lấy mẫu cân bằng, các thực nghiệm toàn diện đã được tiến hành trên bộ dữ liệu CICIDS2017 đã được tinh lọc.

\subsection{Môi trường thực nghiệm}
Khung tính toán được triển khai trên máy trạm Apple Silicon trang bị bộ vi xử lý Apple M3 và 16GB bộ nhớ thống nhất (Unified Memory) chạy hệ điều hành macOS 15.5. Cấu trúc phần mềm được xây dựng trên ngôn ngữ Python 3.12, sử dụng thư viện \texttt{Scikit-Learn} để xây dựng quy trình dữ liệu (data pipeline), điều phối ensemble và mô hình Random Forest; thư viện \texttt{XGBoost} và \texttt{LightGBM} cho tăng cường gradient (gradient boosting); và thư viện \texttt{Imbalanced-Learn} cho chiến lược lấy mẫu cân bằng SMOTEENN.

\subsection{Các chỉ số đánh giá}
Trong các kịch bản phát hiện xâm nhập bị mất cân bằng dữ liệu nghiêm trọng, chỉ số Độ chính xác (Accuracy) tổng thể đơn thuần là một chỉ số gây hiểu lầm. Do đó, đánh giá của chúng tôi sử dụng một bộ các chỉ số toàn diện được rút ra từ các thành phần của ma trận nhầm lẫn (confusion matrix): Dương tính thật ($TP$), Âm tính thật ($TN$), Dương tính giả ($FP$), và Âm tính giả ($FN$).

\begin{itemize}
    \item \textbf{Độ chính xác (Accuracy):} Tỷ lệ tổng thể của các dự đoán đúng trên tổng số quan sát.
    \begin{equation}
        \text{Accuracy} = \frac{TP + TN}{TP + TN + FP + FN}
    \end{equation}

    \item \textbf{Độ chính xác (Precision/Độ xác thực):} Tỷ lệ các dự đoán dương tính đúng trên tổng số các dự đoán là dương tính.
    \begin{equation}
        \text{Precision} = \frac{TP}{TP + FP}
    \end{equation}

    \item \textbf{Độ triệu hồi (Recall/Tỷ lệ phát hiện):} Tỷ lệ các cuộc tấn công được phát hiện đúng trên tổng số các cuộc tấn công thực tế.
    \begin{equation}
        \text{Recall} = \frac{TP}{TP + FN}
    \end{equation}

    \item \textbf{Điểm F1 (F1-Score):} Trung bình điều hòa của Precision và Recall.
    \begin{equation}
        \text{F1} = 2 \times \frac{\text{Precision} \cdot \text{Recall}}{\text{Precision} + \text{Recall}}
    \end{equation}

    \item \textbf{Tỷ lệ Dương tính giả (False Positive Rate - FPR):} Xác suất kích hoạt cảnh báo sai đối với lưu lượng truy cập lành tính. Việc cực tiểu hóa FPR là tối quan trọng để ngăn chặn sự mệt mỏi của các nhà phân tích hệ thống.
    \begin{equation}
        \text{FPR} = \frac{FP}{FP + TN}
    \end{equation}

    \item \textbf{AUC-ROC:} Diện tích dưới đường cong đặc tính hoạt động của bộ thu (Area Under the Receiver Operating Characteristic curve) định lượng khả năng phân biệt giữa các lớp của mô hình trên tất cả các ngưỡng phân loại. Chỉ số AUC-ROC bằng 1,0 cho thấy khả năng phân biệt hoàn hảo.

    \item \textbf{AUC-PR (Độ chính xác trung bình):} Diện tích dưới đường cong Precision-Recall đặc biệt hữu ích cho các bộ dữ liệu mất cân bằng, vì nó tập trung vào hiệu năng của lớp thiểu số mà không bị ảnh hưởng bởi số lượng lớn các trường hợp âm tính thật.
    
    \item \textbf{Độ trễ suy luận (Inference Latency):} Để đánh giá khả năng sẵn sàng triển khai tại biên, chúng tôi đo lường thời gian phân loại đầu cuối cần thiết cho mỗi 1.000 mẫu luồng mạng, được biểu diễn bằng mili giây (ms/1k).
\end{itemize}

\subsection{Giao thức thực nghiệm}
Để đảm bảo tính nghiêm túc khoa học, chúng tôi áp dụng giao thức đánh giá đa khía cạnh:
\begin{enumerate}
    \item \textbf{So sánh với mô hình cơ sở (Baseline):} Mỗi bộ phân loại cơ sở (XGBoost, LightGBM, Random Forest) được huấn luyện và đánh giá riêng lẻ trên cùng một bộ dữ liệu cân bằng, cung cấp sự so sánh trực tiếp để chứng minh giá trị gia tăng của chiến lược ensemble.
    \item \textbf{Nghiên cứu cắt bỏ (Ablation Study):} Chúng tôi loại bỏ hoặc thay thế một cách hệ thống các thành phần của khung đề xuất để định lượng đóng góp riêng lẻ của chúng: (a) loại bỏ SMOTEENN, (b) thay thế Stacking bằng Soft Voting, và (c) loại bỏ bộ phân loại cơ sở Random Forest.
    \item \textbf{Kiểm định chéo 5 lớp phân tầng (Stratified 5-Fold Cross-Validation):} Toàn bộ quy trình (SMOTEENN + Stacking) được đánh giá bằng kiểm định chéo 5 lớp phân tầng để báo cáo giá trị trung bình $\pm$ độ lệch chuẩn, chứng minh tính ổn định và khả năng tái lập của kết quả.
\end{enumerate}
